\documentclass{beamer}

\definecolor{links}{HTML}{2A1B81}
\hypersetup{colorlinks,linkcolor=,urlcolor=links}

\mode<presentation> {

% The Beamer class comes with a number of default slide themes
% which change the colors and layouts of slides. Below this is a list
% of all the themes, uncomment each in turn to see what they look like.

%\usetheme{default}
%\usetheme{AnnArbor}
%\usetheme{Antibes}
%\usetheme{Bergen}
%\usetheme{Berkeley}
%\usetheme{Berlin}
%\usetheme{Boadilla}
\usetheme{CambridgeUS}
%\usetheme{Copenhagen}
%\usetheme{Darmstadt}
%\usetheme{Dresden}
%\usetheme{Frankfurt}
%\usetheme{Goettingen}
%\usetheme{Hannover}
%\usetheme{Ilmenau}
%\usetheme{JuanLesPins}
%\usetheme{Luebeck}
%\usetheme{Madrid}
%\usetheme{Malmoe}
%\usetheme{Marburg}
%\usetheme{Montpellier}
%\usetheme{PaloAlto}
%\usetheme{Pittsburgh}
%\usetheme{Rochester}
%\usetheme{Singapore}
%\usetheme{Szeged}
%\usetheme{Warsaw}

% As well as themes, the Beamer class has a number of color themes
% for any slide theme. Uncomment each of these in turn to see how it
% changes the colors of your current slide theme.

%\usecolortheme{albatross}
%\usecolortheme{beaver}
%\usecolortheme{beetle}
%\usecolortheme{crane}
%\usecolortheme{dolphin}
%\usecolortheme{dove}
%\usecolortheme{fly}
%\usecolortheme{lily}
%\usecolortheme{orchid}
%\usecolortheme{rose}
%\usecolortheme{seagull}
\usecolortheme{seahorse}
%\usecolortheme{whale}
%\usecolortheme{wolverine}

%\setbeamertemplate{footline} % To remove the footer line in all slides uncomment this line
%\setbeamertemplate{footline}[page number] % To replace the footer line in all slides with a simple slide count uncomment this line

%\setbeamertemplate{navigation symbols}{} % To remove the navigation symbols from the bottom of all slides uncomment this line
}
%
%-------------------------------------------------------
% INCLUDE PACKAGES
%-------------------------------------------------------
\usepackage[utf8]{inputenc}
\usepackage[english]{babel}
\usepackage[T1]{fontenc}
\usepackage{helvet}
\usepackage{graphicx} % Allows including images
\usepackage{booktabs} % Allows the use of \toprule, \midrule and \bottomrule in tables
\usepackage{siunitx}
\newcommand{\gc}{\degreeCelsius}
\usepackage{tikz}
\usepackage{listings}
% -------------------------------------------------------
% INCLUDE NOTES
%-------------------------------------------------------
%\setbeameroption{show notes}


%-------------------------------------------------------
% INFORMATION IN THE TITLE PAGE
%-------------------------------------------------------

\title[IntSciProg]{\textbf{Tutorial of GNU/Linux in the console. Module II}}

\subtitle{\textbf{Introduction to Scientific Programming}}


\newcommand\RBox[1]{%
  \tikz\node[draw,rounded corners,align=center,] {#1};%
}  
\author[Curro Pérez-Bernal]
{%
  \texorpdfstring{
    % \begin{columns}
    %   \column{.45\linewidth}
      \centering
      \RBox{Curro Pérez-Bernal\\
        \href{mailto:curropb@uhu.es}{curropb@uhu.es}}
    % \end{columns}
    \vspace{0.5cm}
    \begin{columns}
      \column{.75\linewidth}
      \centering
      \RBox{Integrated Sciences Department --
        University of Huelva}
    \end{columns}
  }
  {John Doe \& Jane Doe}
}
% \author[Jos\'e E. Garc\'{\i}a Ramos]{Jos\'e E. Garc\'{\i}a Ramos}% Your name
% \institute[DCI-UHU]
  
%   \medskip
%   \textit{curropb@uhu.es} % Your email address
  
% \author[Curro P\'erez-Bernal ]{Curro P\'erez Bernal}% Your name
% \institute[DCI-UHU]
  
%   \medskip
%   \textit{curropb@uhu.es} % Your email address

% there must be an empty line above this line - otherwise some unwanted space is added between the university and the country (I do not know why;( )

\date{\today} % Date, can be changed to a custom date

\titlegraphic{\begin{center}\vspace{-0.5cm}\hspace{0.5cm}\includegraphics[width=0.75\linewidth]{Figs/uhu_ceafmc_fisum_isp.png}\end{center}}


%-------------------------------------------------------
% THE BODY OF THE PRESENTATION
%-------------------------------------------------------

\begin{document}

%-------------------------------------------------------
% THE TITLEPAGE
%-------------------------------------------------------

%{\1% % this is the name of the PDF file for the background
\begin{frame}[plain,noframenumbering] % the plain option removes the header from the title page, noframenumbering removes the numbering of this frame only
  \titlepage % call the title page information from above
\end{frame}%}


\begin{frame}{Chapter Outline}{}
\tableofcontents
\end{frame}

\section{Creating Folders}
%%%%%%%%%%%%%%%%%%%%%%%%%%%%%%%%%%%%%%%%%%%%%%%%%%%%%%%%%%%%%%%%%%%%%%%%%%%% 
%%%%%%%%%%%%%%%%%%%%%%%%%%%%%%%%%%%%%%%%%%%%%%%%%%%%%%%%%%%%%%%%%%%%%%%%%%%% 
\begin{frame}[fragile]{Creating directories with \alert{\texttt{mkdir}}}
  % -------------------------------------------------------
  
  \begin{block}{ Syntax: \alert{\texttt{mkdir}} \emph{folder\_name ...} }
    
  {\scriptsize
      \begin{lstlisting}
alice@platea:~$ ls
README  Desktop    Documents  Downloads  
alice@platea:~$
alice@platea:~$ mkdir tmp
alice@platea:~$ ls
README  Desktop    Documents  Downloads  tmp
alice@platea:~$ cd tmp
alice@platea:~/tmp$ 
      \end{lstlisting}
  }
    
    \begin{columns}[t]
      \column{.44\textwidth}
      {\footnotesize \emph{directory\_name} can be
        given as an \alert{absolute} or a \alert{relative} path. The following three commands are
        equivalent.}
      \column{.56\textwidth}
        \hspace{-3cm}
      {\scriptsize
        \begin{lstlisting}
alice@platea:~$ mkdir tmp
alice@platea:~$ mkdir /home/alice/tmp
alice@platea:~$ mkdir ./tmp
alice@platea:~$ mkdir ../alice/tmp
        \end{lstlisting}
      }
    \end{columns}

  \end{block}
  
  
  \note{
    {\tiny
      
      Notes Module I
    }
  }
\end{frame}
%%%%%%%%%%%%%%%%%%%%%%%%%%%%%%%%%%%%%%%%%%%%%%%%%%%%%%%%%%%%%%%%%%%%%%%%%%%%%%%%%%%%%%%%%%%
%%%%%%%%%%%%%%%%%%%%%%%%%%%%%%%%%%%%%%%%%%%%%%%%%%%%%%%%%%%%%%%%%%%%%%%%%%%%%%%%%%%%%%%%%%%
%%%%%%%%%%%%%%%%%%%%%%%%%%%%%%%%%%%%%%%%%%%%%%%%%%%%%%%%%%%%%%%%%%%%%%%%%%%% 
%%%%%%%%%%%%%%%%%%%%%%%%%%%%%%%%%%%%%%%%%%%%%%%%%%%%%%%%%%%%%%%%%%%%%%%%%%%%
\section{Copying Files or Folders}
%%%%%%%%%%%%%%%%%%%%%%%%%%%%%%%%%%%%%%%%%%%%%%%%%%%%%%%%%%%%%%%%%%%%%%%%%%%%
%%%%%%%%%%%%%%%%%%%%%%%%%%%%%%%%%%%%%%%%%%%%%%%%%%%%%%%%%%%%%%%%%%%%%%%%%%%%
\begin{frame}[fragile]{Copying files with \alert{\texttt{cp}}}
  % ------------------------------------------------------


  \begin{block}{The  \texttt{cp} command copies files and directories}
    {\footnotesize
      There are three possible ways of using this command:
    

      \begin{lstlisting}
cp [OPTION]... SOURCE DEST
cp [OPTION]... SOURCE... DIRECTORY
cp [OPTION]... -t DIRECTORY SOURCE...
      \end{lstlisting}
 

  Usage Examples
}

  {\scriptsize

  \begin{columns}[T]
      \column{.44\textwidth}Copy file \texttt{README} to
      \texttt{README.txt} in the same directory.
      \column{.56\textwidth}
        \hspace{-3cm}
%      {\scriptsize
        \begin{lstlisting}
alice@platea:~$ cp README README.txt
alice@platea:~$ 
        \end{lstlisting}
%      }
    \end{columns}
  \begin{columns}[T]
      \column{.44\textwidth} Copy files \texttt{README} and
      \texttt{README.2} to the user's home directory.
      \column{.56\textwidth}
        \hspace{-3cm}
%      {\scriptsize
        \begin{lstlisting}
alice@platea:/tmp$ cp README README.2 ~
alice@platea:/tmp$ 
        \end{lstlisting}
%      }
    \end{columns}
  \begin{columns}[T]
      \column{.44\textwidth} To copy a directory the option
    \texttt{-r} (recursively) should be used. The example shows how to
    copy the directory
    \texttt{\textasciitilde/Images} to the directory \texttt{./tmp}.
      \column{.56\textwidth}
        \hspace{-3cm}
%      {\scriptsize
        \begin{lstlisting}
alice@platea:~$ cp -r ~/Images ./tmp
alice@platea:~$
        \end{lstlisting}
%      }
    \end{columns}
}
  \end{block}
  
%%%%%%%%%%%%%%%%%%%%%%%%%%%%%%%%%%%%%%%%%%%%%%%%%%%%%%%%%%%%%%%%%%%%%%%%%%%%%%%%%%
\note{
{\tiny

Notes Module I
}
}
\end{frame}
%%%%%%%%%%%%%%%%%%%%%%%%%%%%%%%%%%%%%%%%%%%%%%%%%%%%%%%%%%%%%%%%%%%%%%%%%%%%
%%%%%%%%%%%%%%%%%%%%%%%%%%%%%%%%%%%%%%%%%%%%%%%%%%%%%%%%%%%%%%%%%%%%%%%%%%%%
%%%%%%%%%%%%%%%%%%%%%%%%%%%%%%%%%%%%%%%%%%%%%%%%%%%%%%%%%%%%%%%%%%%%%%%%%%%% 
%%%%%%%%%%%%%%%%%%%%%%%%%%%%%%%%%%%%%%%%%%%%%%%%%%%%%%%%%%%%%%%%%%%%%%%%%%%%
\section{Moving Files}
%%%%%%%%%%%%%%%%%%%%%%%%%%%%%%%%%%%%%%%%%%%%%%%%%%%%%%%%%%%%%%%%%%%%%%%%%%%%
%%%%%%%%%%%%%%%%%%%%%%%%%%%%%%%%%%%%%%%%%%%%%%%%%%%%%%%%%%%%%%%%%%%%%%%%%%%%
%%%%%%%%%%%%%%%%%%%%%%%%%%%%%%%%%%%%%%%%%%%%%%%%%%%%%%%%%%%%%%%%%%%%%%%%%%%%
%%%%%%%%%%%%%%%%%%%%%%%%%%%%%%%%%%%%%%%%%%%%%%%%%%%%%%%%%%%%%%%%%%%%%%%%%%%%
\begin{frame}[fragile]{Moving files with \alert{\texttt{mv}}}
  % ------------------------------------------------------


  \begin{block}{The  \texttt{mv} command moves files and directories}
    {\footnotesize
      There are three possible ways of using this command:
    

      \begin{lstlisting}
mv [OPTION]... SOURCE DEST
mv [OPTION]... SOURCE... DIRECTORY
mv [OPTION]... -t DIRECTORY SOURCE...
      \end{lstlisting}
 

  Usage Examples
}

  {\scriptsize

  \begin{columns}[T]
      \column{.44\textwidth} Rename file \texttt{README} to
      \texttt{README.txt} in the same directory.
      \column{.56\textwidth}
        \hspace{-3cm}
%      {\scriptsize
        \begin{lstlisting}
alice@platea:~$ mv README README.txt
alice@platea:~$ 
        \end{lstlisting}
%      }
    \end{columns}
  \begin{columns}[T]
      \column{.44\textwidth} Move files \texttt{README} and
      \texttt{README.2} to the user's home directory.
      \column{.56\textwidth}
        \hspace{-3cm}
%      {\scriptsize
        \begin{lstlisting}
alice@platea:/tmp$ mv README README.2 ~
alice@platea:/tmp$ 
        \end{lstlisting}
%      }
    \end{columns}
  \begin{columns}[T]
      \column{.44\textwidth} Move directory
    \texttt{\textasciitilde/Images} to directory \texttt{./tmp}.
      \column{.56\textwidth}
        \hspace{-3cm}
%      {\scriptsize
        \begin{lstlisting}
alice@platea:~$ mv ~/Images ./tmp
alice@platea:~$
        \end{lstlisting}
%      }
    \end{columns}
}
  \end{block}
  
%%%%%%%%%%%%%%%%%%%%%%%%%%%%%%%%%%%%%%%%%%%%%%%%%%%%%%%%%%%%%%%%%%%%%%%%%%%%%%%%%%
\note{
{\tiny

Notes Module I
}
}
\end{frame}
%%%%%%%%%%%%%%%%%%%%%%%%%%%%%%%%%%%%%%%%%%%%%%%%%%%%%%%%%%%%%%%%%%%%%%%%%%%%
%%%%%%%%%%%%%%%%%%%%%%%%%%%%%%%%%%%%%%%%%%%%%%%%%%%%%%%%%%%%%%%%%%%%%%%%%%%%
%%%%%%%%%%%%%%%%%%%%%%%%%%%%%%%%%%%%%%%%%%%%%%%%%%%%%%%%%%%%%%%%%%%%%%%%%%%% 
%%%%%%%%%%%%%%%%%%%%%%%%%%%%%%%%%%%%%%%%%%%%%%%%%%%%%%%%%%%%%%%%%%%%%%%%%%%%
\section{Deleting Files and Directories}
%%%%%%%%%%%%%%%%%%%%%%%%%%%%%%%%%%%%%%%%%%%%%%%%%%%%%%%%%%%%%%%%%%%%%%%%%%%%
%%%%%%%%%%%%%%%%%%%%%%%%%%%%%%%%%%%%%%%%%%%%%%%%%%%%%%%%%%%%%%%%%%%%%%%%%%%%
%%%%%%%%%%%%%%%%%%%%%%%%%%%%%%%%%%%%%%%%%%%%%%%%%%%%%%%%%%%%%%%%%%%%%%%%%%%%
%%%%%%%%%%%%%%%%%%%%%%%%%%%%%%%%%%%%%%%%%%%%%%%%%%%%%%%%%%%%%%%%%%%%%%%%%%%%
\begin{frame}[t,fragile]{Deleting folders with \alert{\texttt{rmdir}}}
  % ------------------------------------------------------

  \vspace{-0.3cm}
  \begin{block}{The  \texttt{rmdir} command  removes \emph{empty} directories}
    {\footnotesize
      \begin{lstlisting}
rmdir [OPTION]... DIRECTORY...
      \end{lstlisting}
 

  Usage Examples
}


\vspace{-0.3cm}
{\scriptsize

  \begin{columns}
      \column{.44\textwidth}  Create directories \texttt{READ0}
    and  \texttt{READ1} and remove them.
      \column{.56\textwidth}
%      {\scriptsize
        \begin{lstlisting}
alice@platea:~$ mkdir READ0 READ1
alice@platea:~$ ls -d READ0 READ1
READ0  READ1
alice@platea:~$ rmdir READ0 READ1
alice@platea:~$ ls -d READ0 READ1
ls: cannot access READ0: 
         No such file or directory
ls: cannot access READ1: 
          No such file or directory
        \end{lstlisting}
%      }
    \end{columns}
  \begin{columns}
      \column{.44\textwidth}  The \texttt{rmdir} command only
    removes empty folders.
      \column{.56\textwidth}
%      {\scriptsize
        \begin{lstlisting}
alice@platea:~$ mkdir READ0
alice@platea:~$ cp .bash_history READ0 
alice@platea:~$ rmdir READ0
rmdir: failed to remove `READ0/': 
               Directory not empty
alice@platea:~$
        \end{lstlisting}
%      }
    \end{columns}
}
  \end{block}
  
%%%%%%%%%%%%%%%%%%%%%%%%%%%%%%%%%%%%%%%%%%%%%%%%%%%%%%%%%%%%%%%%%%%%%%%%%%%%%%%%%%
\note{
{\tiny

Notes Module I
}
}
\end{frame}
%%%%%%%%%%%%%%%%%%%%%%%%%%%%%%%%%%%%%%%%%%%%%%%%%%%%%%%%%%%%%%%%%%%%%%%%%%%%
%%%%%%%%%%%%%%%%%%%%%%%%%%%%%%%%%%%%%%%%%%%%%%%%%%%%%%%%%%%%%%%%%%%%%%%%%%%%
\begin{frame}[t,fragile]{Deleting files (and folders) with \alert{\texttt{rm}}}
  % ------------------------------------------------------

  \vspace{-0.3cm}
  \begin{block}{The  \texttt{rm} command  and file deletion}
    {\footnotesize
The command \texttt{rm} removes files. By default it does not remove
directories. Be careful, there does not exist something like an
\textbf{undelete} command.

\begin{lstlisting}
rm [OPTION]... FILE...
\end{lstlisting}
 

  Usage Examples
}


\vspace{-0.3cm}
{\scriptsize

  \begin{columns}
      \column{.44\textwidth}  Remove files \texttt{README} and
      \texttt{README.txt}.
      \column{.56\textwidth}
%      {\scriptsize
        \begin{lstlisting}
alice@platea:~$ rm README README.txt
alice@platea:~$ 
        \end{lstlisting}
%      }
    \end{columns}
  \begin{columns}
      \column{.44\textwidth}  Frequently used options.
      \column{.56\textwidth}
%      {\scriptsize
        \begin{lstlisting}
-i   prompt before every removal
-f   ignore nonexistent files, never prompt
-r   remove directories and 
     their contents recursively
        \end{lstlisting}
%      }
    \end{columns}
  }
  {\footnotesize
   It is important to be very
    careful when using \alert{\texttt{rm -r}} as it is easy to wipe out
    important files or even an entire system.\\ \emph{With great power comes
    great responsability}...}
  \end{block}
  
%%%%%%%%%%%%%%%%%%%%%%%%%%%%%%%%%%%%%%%%%%%%%%%%%%%%%%%%%%%%%%%%%%%%%%%%%%%%%%%%%%
\note{
{\tiny

Notes Module I
}
}
\end{frame}
%%%%%%%%%%%%%%%%%%%%%%%%%%%%%%%%%%%%%%%%%%%%%%%%%%%%%%%%%%%%%%%%%%%%%%%%%%%%
%%%%%%%%%%%%%%%%%%%%%%%%%%%%%%%%%%%%%%%%%%%%%%%%%%%%%%%%%%%%%%%%%%%%%%%%%%%%
\begin{frame}{Saving keystrokes...}{\emph{Programmers can be lazy.} Larry Wall}
%-------------------------------------------------------
  \vspace{-0.3cm}
  \begin{block}{Saving on console typing}
    {\footnotesize Several key combinations save on typing, in particular \texttt{bash-completion} ``magic'' key \alert{\texttt{TAB}}.
    
    
    Usage Examples:
  }
  
{  \scriptsize
  \begin{itemize}
\item The \texttt{TAB} key is a very useful help. It autocompletes any command or file name, if there is only one possible option,  or gives a list of possible options. (``magic'' key).\pause
\item The command \texttt{history} provides a list of typed commands.\pause
\item The user can navigate in the \texttt{history} record with the \texttt{up}-arrow (also \texttt{CTRL-p}) and the \texttt{down}-arrow (also \texttt{CTRL-n}).\pause
\item Selecting text, the key combinations \texttt{Ctrl-Insert} or \texttt{Shift-Ctrl-c} (copy) and \texttt{Shift-Insert} or \texttt{Shift-Ctrl-v} (paste).\pause
\item Selecting text in a console and clicking on the second mouse button yanks this text in the prompt.
\end{itemize}

}  
  \end{block}
%%%%%%%%%%%%%%%%%%%%%%%%%%%%%%%%%%%%%%%%%%%%%%%%%%%%%%%%%%%%%%%%%%%%%%%%%%%%%%%%%%
\note{
{\tiny

Notes Module I
}
}
\end{frame}
%%%%%%%%%%%%%%%%%%%%%%%%%%%%%%%%%%%%%%%%%%%%%%%%%%%%%%%%%%%%%%%%%%%%%%%%%%%%
%%%%%%%%%%%%%%%%%%%%%%%%%%%%%%%%%%%%%%%%%%%%%%%%%%%%%%%%%%%%%%%%%%%%%%%%%%%%
\begin{frame}[t,fragile]{Using \alert{globbing}}
  % ------------------------------------------------------

  \vspace{-0.3cm}
  \begin{block}{File wildcards}
    {\footnotesize
File wildcards save typing, affecting several files simultaneously. The character \alert{\texttt{*}} matches any string, of any length.

\begin{tabular}{|l|l|}
\hline
\texttt{*} & Matches any string, of any length \\
\texttt{bla*} & Matches any string beginning with bla \\
\texttt{*x*} & Matches any string with an \texttt{x} at any position. \\
\texttt{*dat} & Matches any string finishing with \texttt{dat} \\
\hline
\end{tabular}

  Usage Examples
}


\vspace{-0.3cm}
{\scriptsize

  \begin{columns}
      \column{.44\textwidth} Move files ending in \texttt{.txt} in the working directory to the \texttt{/tmp} directory.
      \column{.56\textwidth}
%      {\scriptsize
        \begin{lstlisting}
alice@platea:~$ mv ./*.txt /tmp
alice@platea:~$ 
        \end{lstlisting}
%      }
    \end{columns}
  \begin{columns}
      \column{.44\textwidth}   Copy files beginning with \texttt{output} and ending with \texttt{.dat} to the user's home subdirectory \texttt{Results}.
      \column{.56\textwidth}
%      {\scriptsize
        \begin{lstlisting}
alice@platea:~$ cp output*.dat ~/Results
alice@platea:~$
        \end{lstlisting}
%      }
    \end{columns}
  \begin{columns}
      \column{.44\textwidth}  Delete all files in the current directory beginning with \texttt{a} and ending with \texttt{.ps}.
      \column{.56\textwidth}
%      {\scriptsize
        \begin{lstlisting}
alice@platea:~$ rm ./a*.ps
alice@platea:~$
        \end{lstlisting}
%      }
    \end{columns}
  }
  \end{block}
  
%%%%%%%%%%%%%%%%%%%%%%%%%%%%%%%%%%%%%%%%%%%%%%%%%%%%%%%%%%%%%%%%%%%%%%%%%%%%%%%%%%
\note{
{\tiny

Notes Module I
}
}
\end{frame}
%%%%%%%%%%%%%%%%%%%%%%%%%%%%%%%%%%%%%%%%%%%%%%%%%%%%%%%%%%%%%%%%%%%%%%%%%%%%
%%%%%%%%%%%%%%%%%%%%%%%%%%%%%%%%%%%%%%%%%%%%%%%%%%%%%%%%%%%%%%%%%%%%%%%%%%%%
\begin{frame}[t,fragile]{Using spaces and special characters in file names}
  % ------------------------------------------------------

  \vspace{-0.3cm}
  \begin{block}{Spaces and special characters}
    {\footnotesize
The use of spaces in filenames is allowed, though it is discouraged. Special characters use is even more controversial. Take into account that \alert{space} is the separator between arguments in \texttt{UNIX}.

Files differing on the use of upper and lowercase letters are \alert{different}.



  Usage Examples
}


%\vspace{-0.3cm}
{\scriptsize

  \begin{columns}
      \column{.3\textwidth}Remove a file called \texttt{``my notes.txt''}.
      \column{.7\textwidth}
%      {\scriptsize
        \begin{lstlisting}
alice@platea:~$ rm my\ notes.txt
alice@platea:~$ 
        \end{lstlisting}
%      }
    \end{columns}
  \begin{columns}
      \column{.3\textwidth}    Removing a file with spaces. Check the unexpected result.
      \column{.7\textwidth}
%      {\scriptsize
% alice@platea:~$ ls 
% alice@platea:~$ file.txt text.txt
      \begin{lstlisting}
alice@platea:~$ rm this is a text.txt
rm: cannot remove `this': No such file or directory
rm: cannot remove `is': No such file or directory
rm: cannot remove `a': No such file or directory
alice@platea:~$ rm ``this is a text.txt''
                OR
alice@platea:~$ rm this\ is\ a\ text.txt
alice@platea:~$
        \end{lstlisting}
%      }
    \end{columns}
  }
  \end{block}
  
%%%%%%%%%%%%%%%%%%%%%%%%%%%%%%%%%%%%%%%%%%%%%%%%%%%%%%%%%%%%%%%%%%%%%%%%%%%%%%%%%%
\note{
{\tiny

Notes Module I
}
}
\end{frame}
%%%%%%%%%%%%%%%%%%%%%%%%%%%%%%%%%%%%%%%%%%%%%%%%%%%%%%%%%%%%%%%%%%%%%%%%%%%%
%%%%%%%%%%%%%%%%%%%%%%%%%%%%%%%%%%%%%%%%%%%%%%%%%%%%%%%%%%%%%%%%%%%%%%%%%%%%


%%%%%%%%%
\end{document}
